% Code areas by their number of compounds, on half-decade bins, by the digits of their name range (registry/areas.csv.gz).
% Data: area-sizes.dat, written by tools/paper-figures.py.
\documentclass[tikz,border=3pt]{standalone}
\usepackage{fontspec}
\setmainfont{LibertinusSans}[Extension=.otf,UprightFont=*-Regular,BoldFont=*-Bold,ItalicFont=*-Italic]
\usepackage{unicode-math}
\setmathfont{LibertinusMath-Regular.otf}
\usepackage{pgfplots}
\pgfplotsset{compat=1.18}
\usepackage{pgfplotstable}
% English by default; the French paper compiles <name>-fr.tex, which sets \fr to 1 and inputs this file.
\providecommand{\fr}{0}
\newcommand{\T}[2]{\ifnum\fr=1 #2\else #1\fi}
\ifnum\fr=1 \pgfkeys{/pgf/number format/.cd, use comma, 1000 sep={\,}}\else \pgfkeys{/pgf/number format/.cd, 1000 sep={,}}\fi
\definecolor{ink}{HTML}{00664F}\definecolor{flaggreen}{HTML}{007A5E}\definecolor{flagred}{HTML}{CE1126}
\definecolor{flagyellow}{HTML}{FCD116}\definecolor{amber}{HTML}{C98A00}
\begin{document}
\begin{tikzpicture}
\begin{axis}[ybar stacked, width=12cm, height=6.6cm, bar width=17pt, ymin=0, ymax=7800,
  xtick={0,...,8}, xticklabels={1–2,3–9,10–31,32–99,100–315,316–999,\T{1,000–3,161}{1\,000–3\,161},\T{3,162–9,999}{3\,162–9\,999},\T{10,000–19,107}{10\,000–19\,107}},
  x tick label style={font=\scriptsize, rotate=25, anchor=north east}, y tick label style={font=\footnotesize},
  xlabel={\T{compounds in the area}{concessions dans la zone}}, ylabel={\T{areas}{zones}}, label style={font=\small},
  axis line style={gray}, ymajorgrids, grid style={gray!15}, xmin=-0.6, xmax=8.6,
  legend style={at={(0.02,0.97)}, anchor=north west, font=\footnotesize, draw=gray!50}, legend cell align=left,
  scaled y ticks=false]
\addplot[fill=flaggreen!55, draw=none] table[x=bin, y=d3] {area-sizes.dat};
\addlegendentry{\T{3 digits}{3 chiffres}}
\addplot[fill=flaggreen, draw=none] table[x=bin, y=d4] {area-sizes.dat};
\addlegendentry{\T{4 digits}{4 chiffres}}
\addplot[fill=flagred, draw=none] table[x=bin, y=d5] {area-sizes.dat};
\addlegendentry{\T{5 digits}{5 chiffres}}
\end{axis}
\end{tikzpicture}
\end{document}
